% Generated draft: requires mathematical and source review.
\documentclass[11pt]{article}
\usepackage[margin=25mm]{geometry}
\usepackage{fontspec}
\setmainfont{Latin Modern Roman}
\usepackage{amsmath,amssymb,mathtools,unicode-math}
\setmathfont{Latin Modern Math}
\usepackage{enumitem,needspace,xurl,hyperref}
\setlength{\parindent}{0pt}
\setlength{\parskip}{5pt}
\setlength{\emergencystretch}{3em}
\newcommand{\N}{\mathbb{N}}
\newcommand{\Z}{\mathbb{Z}}
\newcommand{\Q}{\mathbb{Q}}
\newcommand{\R}{\mathbb{R}}
\newcommand{\C}{\mathbb{C}}
\newcommand{\F}{\mathbb{F}}
\newcommand{\A}{\mathbb{A}}
\newcommand{\PP}{\mathbb{P}}
\newcommand{\E}{\mathbb{E}}
\newcommand{\eps}{\varepsilon}
\newcommand{\dd}{\,\mathrm{d}}
\begin{document}
\section*{UnsolvedMath ID: 800043}
\textbf{Source problem code:} catalogue entry
\subsection{Definitions and mathematical statement}
A semigroup is a set equipped with an associative binary operation, written multiplicatively. An idempotent semigroup, also called a band, satisfies
\[
\forall x\in S,\qquad x^2=x.
\]
A semigroup variety is a class of semigroups axiomatized by identities. Let \(\mathcal B\) denote the variety of all bands, defined by associativity and the identity \(x^2=x\).

\textbf{Source question (G. Mashevitzky).} Does there exist a semigroup variety \(\mathcal V\subseteq\mathcal B\) that is both \emph{inclusive} and \emph{nonelementary}, in the senses intended by the source?

The supplied record does not define either ``inclusive'' or ``nonelementary.'' These are unresolved qualifications on \(\mathcal V\); no interpretation of them, including a model-theoretic interpretation of ``elementary,'' is assumed here. Consequently, the supplied evidence does not determine a fully specified mathematical existence statement.
\subsection{Short English statement}
G. Mashevitzky asks whether there is an inclusive, nonelementary variety of idempotent semigroups. The two qualifying terms require source definitions before the question can be made fully precise. The local record labels the problem open; this is supplied status metadata, not an independently verified assessment.
\subsection{Sources}
\begin{itemize}
\item UnsolvedMath local registry, permanent ID \(800043\), exact source problem code \texttt{catalogue entry}, title ``The Lyapin's Notebook --- Problem 45.'' The supplied record attributes the question itself to G. Mashevitzky; its proposer metadata lists Bershadsky, Kublanovsky, and Mashevitzky and gives the year 2026.
\item Bershadsky, Kublanovsky, Mashevitzky, \emph{The Lyapin's Notebook: Unsolved Problems in Semigroup Theory} (2026), Problem 45. Supplied source URL: \url{https://arxiv.org/abs/2604.04763}. The record reports access on 2026-10-16 and extraction by \texttt{arxiv\_latex}. These are supplied provenance details; the source has not been independently consulted or verified for this draft.
\end{itemize}
\end{document}
